\scanpage{601}

\mspage{10}

En effet, la donnée d'une telle famille $(E_{\unsure{\xi|\zeta}})_{\xi\in k}$ équivaut à la donnée d'une application $E\unsure{-|\to}k$, i.e. d'un \unsure{élmt} de $\unsure{f_E|k_E}$. \deleted{\unsure{Et}}

\deleted{\unsure{S|B}} \uline{Scholie} La propriété qu'on vient d'expliciter, pour un filtre $\Phi$\inserted{, sur $E$}, ne dépend \deleted{que} de $k$ que par l'intermédiaire de $\unsure{\pi|\kappa}=\operatorname{card}k$. On peut aussi l'énoncer en disant que pour tte \uline{partition} $(E_i)_{i\in I}$ de $E$ en \inserted{un ens. $I$ de} parties tel que $\operatorname{card}I\leq\pi$, il existe un (et un seul) des $E_i$ qui soit dans $\Phi$. On dit alors que $\unsure{Q|\Phi}$ est un $\pi$-ultrafiltre.

Si $\pi\geq 2$ est fini, cette propriété équivaut à dire que $\Phi$ est un ultrafiltre. \deleted{Si $\pi$ est infini, il n'est pas clair qu'il \unsure{existe|a}} \inserted{Par ailleurs, quel que soit $\pi$, les} ultrafiltres $\Phi$ \uline{triviaux} \unsure{(}(i.e. tels qu'il existe $\xi\in E$ avec $\{\xi\}\in\Phi$, ou encore \deleted{\unsure{$\Phi$}} tel que $\Phi=\{A\in\mathcal P(E)\mid\xi\in A\}$) sont des $\pi$-ultrafiltres. \deleted{Mais} Si $\pi$ est transfini, il n'est pas clair qu'il existe des $\pi$-ultrafiltres non triviaux. On sait qu'il n'en existe sur $E$ \uline{que} si \deleted{\unsure{card}} $E\notin\unsure{\mathcal U|U}_1$, où $\mathcal U_1$ est le plus petit univers \deleted{tel que \unsure{l'ensemble}} tel que $\mathbb N\in\mathcal U_1$, i.e tel

% archive sheet number: 600
